% Additive continuation. Do not overwrite the accepted SGT/NCT/TAC sources.
\begingroup
\providecommand{\rank}{\operatorname{rank}}
\section{A polynomial forward bound for the full weighted conductor}
\label{sec:WCF}

The original Split-Zero programme and the deposited source notes retain
their existing attribution \cite{human:globalization2026,human:splitzero2026}.
The new calculation in this section acts on the literal SGT10 matrix. It
improves its forward bound at ranks of order $a^2$. Its exact triangular
determinant then controls every inverse exterior rank and supplies a
two-sided four-endpoint comparison for the entire conductor block.

\subsection{Original spaces, mass, phases and matrix}

Fix the original hypothetical simple quartet, its original period and
logarithm branch. Retain $0<\delta<1/2$, $\gamma>2$, and all original
coefficients $a_{uw}$, including repeated contributions. Put
\[
\begin{gathered}
 a=4l+1\ge9,\quad q=(a+1)^2,\quad q'=(a-7)^2,\quad
 \beta_{n,u,w}=n/2+(2u-n)\delta+i(2w-n)\gamma,\\
 E_A(z)=\sum_{u,w=0}^{8}a_{uw}e^{\beta_{8,u,w}z},\quad
 \mu_n=E_A^{(n)}(0),\quad v=\operatorname{ord}_0E_A\le80,
 \quad \mu_v\ne0,\quad A_{\rm abs}=\sum_{u,w=0}^{8}|a_{uw}|,\\
 s\in\{1,a\},\quad b=s/2,\quad c_0=a/2,\quad c_1=a/2-4,
 \quad N=D+v,\quad D\ge0.
\end{gathered}\tag{WCF1}
\]
The symbol and its coefficients are fixed as $a$ varies. The Hermitian
inner product is conjugate-linear in its first variable. The original
polynomial norm, orthonormal frame, and leading coefficient are
\[
\begin{gathered}
 dm_s(y)=\frac{(2\pi)^{s/2}2^{s/2}}{4\pi\Gamma(s/2)}
       |\Gamma(s/4+iy/2)|^2\,dy,\quad M_s=(2\pi)^{s/2},\\
 \|P\|_{s,c}^2=\int_{\mathbb R}|P(c+iy)|^2dm_s(y),\qquad
 p_{0,s}=1,\quad p_{1,s}=y,\\
 p_{n+1,s}=yp_{n,s}-n(n-1+s/2)p_{n-1,s},\quad
 h_{n,s}=M_sn!(s/2)_n,\\
 \phi_{n,s,c}(S)=\frac{p_{n,s}((S-c)/i)}{\sqrt{h_{n,s}}},\qquad
 [S^n]\phi_{n,s,c}=\frac{i^{-n}}{\sqrt{h_{n,s}}},\qquad
 \rho_n=\sqrt{\frac{n!}{(s/2)_n}}.
\end{gathered}\tag{WCF2}
\]
The classical dictionary is
$p_{n,s}(y)=n!P_n^{(s/4)}(y/2;\pi/2)$, with $dy=2\,d(y/2)$.
The weight, norm, recurrence and generating function are respectively
DLMF (18.19.8)--(18.19.9), (18.22.8), and (18.23.7)
\cite{human:dlmf18}. In the generating function take
$\lambda=s/4>0$, $x=y/2$, $\varphi=\pi/2$, and $z=t$, $|t|<1$.
The two powers take their branches with value one at $t=0$; their product is
\[
 \sum_{n\ge0}\frac{p_{n,s}(y)}{n!}t^n
   =(1+t^2)^{-s/4}e^{y\arctan t}.
 \tag{WCF3}
\]
This also follows by multiplying WCF2's recurrence by $t^n/n!$ and
solving $(1+t^2)\partial_t F=(y-(s/2)t)F$, $F(y,0)=1$.
Thus the stated generating identity holds at the complex polynomial
arguments below, without a change of phase or of the measure.

For $\mathcal P_d=\mathbb C[S]_{\le d}$, $\mathcal P_{-1}=0$, the
unscaled conductor map is
\[
 \mathcal T_A P(S')=\sum_{u,w=0}^{8}a_{uw}P(S'+\beta_{8,u,w}),\qquad
 A_{D,s}:(\mathcal P_{D+v}/\mathcal P_{v-1},\|\cdot\|_{s,c_0,\rm quot})
       \longrightarrow(\mathbb C[S']_{\le D},\|\cdot\|_{s,c_1}).
 \tag{WCF4}
\]
Indeed $\mathcal T_A S^n=\sum_{j=v}^n\binom nj\mu_j S'^{n-j}$.
For $n\ge v$ the leading coefficient is $\binom nv\mu_v\ne0$;
for $n<v$ the image is zero. Successive elimination in degree proves
that WCF4 is an isomorphism with precisely the displayed kernel.
Its quotient frame is represented by
$\phi_{v,s,c_0},\ldots,\phi_{D+v,s,c_0}$, the orthogonal complement
of that full kernel. Set
\[
 w(t)=-i\arctan t=\tfrac12\log\frac{1-it}{1+it},\qquad
 g(t)=t^{-v}e^{-4w(t)}E_A(w(t)),\quad
 g(0)=(-i)^v\mu_v/v!,\quad g(t)=\sum_{j\ge0}g_jt^j.
 \tag{WCF5}
\]
The singularity at zero is removable. The displayed logarithm is the
analytic branch equal to zero there. Apply the finite sum WCF4 to WCF3.
Since $c_0-c_1=4$, the result is $t^vg(t)$ times the generating
function centred at $c_1$. Comparing coefficients, including
$n!/\sqrt{h_{n,s}}=\rho_n/\sqrt{M_s}$, proves
\[
 A_{D,s}=\operatorname{diag}(\rho_0,\ldots,\rho_D)^{-1}
          T_D(g)\operatorname{diag}(\rho_v,\ldots,\rho_{D+v}),
 \quad
 (A_{D,s})_{ij}=\begin{cases}g_{j-i}\rho_{j+v}/\rho_i,&i\le j,\\0,&i>j.
 \end{cases}
 \tag{WCF6}
\]
Here $T_D(g)_{ij}=g_{j-i}$ above the diagonal. The two copies of
$M_s^{-1/2}$ cancel in this matrix because the source order is the same
on both sides. All norms and all endpoint Grams retain their original
mass. No scalar $v!/\mu_v$ is inserted in WCF4 or WCF6.

\subsection{Forward coefficient and weighted-shift estimates}

For $|t|\le r<1$ the Cayley ratio $z=(1-it)/(1+it)$ satisfies
\[
 \Re z=\frac{1-|t|^2}{|1+it|^2}>0,\qquad
 \frac{1-r}{1+r}\le|z|\le\frac{1+r}{1-r},\qquad
 |\arg z|<\pi/2.
 \tag{WCF7}
\]
The first identity follows by multiplying by $1-i\bar t$; the
modulus bounds follow from the triangle inequality and its reverse.
The branch in WCF5 is therefore exactly the logarithm on the right
half-plane. For $\beta_{8,u,w}-4=x+iy$, we have
$|x|\le8\delta$, $|y|\le8\gamma$. Consequently
\[
 \left|e^{(\beta_{8,u,w}-4)w(t)}\right|
 =\exp\!\left(\tfrac{x}{2}\log|z|-\tfrac{y}{2}\arg z\right)
 \le\left(\frac{1+r}{1-r}\right)^{4\delta}e^{2\pi\gamma}.
 \tag{WCF8}
\]
This estimate retains the imaginary shifts through $e^{2\pi\gamma}$;
it does not replace them by a real-root symbol.

For $D\ge1$, choose $r_D=D/(D+1)$ and define
\[
 C_D=e\,2^v A_{\rm abs}e^{2\pi\gamma}(2D+1)^{4\delta}.
 \tag{WCF9}
\]
On $|t|=r_D$, WCF5 and WCF8 bound $|g(t)|$ by
$r_D^{-v}A_{\rm abs}e^{2\pi\gamma}(2D+1)^{4\delta}$.
Cauchy's coefficient integral gives $|g_l|\le r_D^{-l}\max|g|$.
For $0\le l\le D$, $r_D^{-l}\le(1+1/D)^D\le e$ and
$r_D^{-v}\le2^v$. Hence
\[
 |g_l|\le C_D\qquad(0\le l\le D).
 \tag{WCF10}
\]

The $l$th upper diagonal in WCF6 is a weighted shift whose norm is
the maximum of $\rho_{i+l+v}/\rho_i$. Its nonzero columns have
orthogonal images, which proves this norm assertion directly. If
$s=a\ge9$, then $b>1$ and $\rho_{n+1}/\rho_n=
\sqrt{(n+1)/(n+b)}\le1$, so this norm is at most one.
If $s=1$, the factors in
\[
 \frac{\rho_{i+d}^2}{\rho_i^2}
 =\prod_{t=0}^{d-1}\left(1+\frac{1/2}{i+t+1/2}\right)
 \le\frac{d!}{(1/2)_d}\le2d+1
 \qquad(d\ge0)
 \tag{WCF11}
\]
decrease with $i$. For the last bound use induction: it is equality
at $d=0$, and multiplication by $(d+1)/(d+1/2)$ takes $2d+1$
to $2d+2\le2d+3$. The following integral calculation avoids paying
for the number of upper diagonals.

For $s=a$, so $b>1$, give a polynomial $P(z)=\sum p_nz^n$ the norm
\[
 \|P\|_{\rho}^2=\sum_n|p_n|^2\rho_n^2
  =(b-1)\int_0^1\int_0^{2\pi}|P(\sqrt{x}e^{i\theta})|^2
       \frac{d\theta}{2\pi}(1-x)^{b-2}\,dx.
 \tag{WCF11a}
\]
Angular integration kills unequal monomials. The radial identity
$(b-1)\int_0^1x^n(1-x)^{b-2}dx=n!/(b)_n$ follows either by
integration by parts and its value at $n=0$, or directly from Euler's
beta integral, DLMF (5.12.1), at the positive parameters $n+1,b-1$
\cite{human:dlmf5}. Thus WCF11a is an isometric coefficient
realization of these exact positive weights. It is an auxiliary proof
space for WCF6; the original Gamma measure WCF2 is unchanged.

Multiplication by $g(rz)$, $0<r<1$, has norm at most
$\max_{|t|=r}|g(t)|$ in WCF11a, since that pointwise modulus bound
holds on the whole unit disk. Its power series converges uniformly
on the closed unit disk. Orthogonal projection onto degrees at most
$D$ contracts the norm. In the orthonormal basis $z^n/\rho_n$, the
transpose of this compressed multiplication matrix is
\[
 U_r=\operatorname{diag}(\rho_0,\ldots,\rho_D)^{-1}
            T_D(g(r\,\cdot))\operatorname{diag}(\rho_0,\ldots,\rho_D).
 \tag{WCF11b}
\]
Transposition preserves operator norm. For
$R_r=\operatorname{diag}(1,r,\ldots,r^D)$, the unscaled matrix is
$U=R_rU_rR_r^{-1}$. Hence
$\|U\|\le r^{-D}\max_{|t|=r}|g(t)|$. Finally
$A_{D,a}=U\operatorname{diag}(\rho_{i+v}/\rho_i)$, whose last
diagonal is a contraction. Taking $r=r_D$ and WCF8 proves
$\|A_{D,a}\|\le C_D$.

For $s=1$ apply the same angular integration with the unweighted
circle norm $\sum|p_n|^2$: multiplication and orthogonal compression
give $\|T_D(g)\|\le r^{-D}\max_{|t|=r}|g(t)|$ by the same radial
similarity. In this case $\rho_n$ increases, so the two diagonal
operator norms in WCF6 have product $\rho_{D+v}$, bounded by
$\sqrt{2(D+v)+1}$ in WCF11. We have therefore proved the sharper
finite estimate
\[
 \boxed{\quad \|A_{D,s}\|\le B^+_{D,s}:=
 \begin{cases}
 C_D\sqrt{2(D+v)+1},&D\ge1,\ s=1,\\
 C_D,&D\ge1,\ s=a,\\
 |\mu_v|\rho_v/v!,&D=0.
 \end{cases}\quad}
 \tag{WCF12}
\]
The $D=0$ case is the exact scalar modulus from WCF6; $\rho_0=1$.
The estimate therefore includes $v=0$ and the smallest cutoff.
In particular, uniformly for $s\in\{1,a\}$ and $D\le C a^2$,
\[
 \log^+B^+_{D,s}=O_{\rm actual,C}(\log(a+1)).
 \tag{WCF13}
\]
For $D\ge1$ this follows by taking the logarithm of WCF12 and
using fixed $v,\delta,\gamma,A_{\rm abs}$. At $D=0$, $\rho_v\le1$
for $s=a$ and is fixed for $s=1$. This proves the asserted uniformity.

\subsection{A finite inverse estimate with every principal part}

The following records the inverse terms needed at the negative endpoints.
It also prevents an upper-triangular estimate from being applied to a
lower-triangular correction. Retain the original SGT zero-free constants
\[
\begin{gathered}
 B_0=\max\{1,\max_{u,w}|\beta_{8,u,w}|\},\quad
 C_v=A_{\rm abs}B_0^v/|\mu_v|,\\
 R_0=B_0^{-1}\min\{1,(v+1)/(2eC_v)\},\quad
 a_0=\min\{1/2,\tanh(R_0/2)\},\\
 d_0=1-\frac{a_0^2}{3(1-a_0^2)},\qquad
 M_0=\frac{2v!}{|\mu_v|}d_0^{-v}e^{2R_0},\qquad
 S_D(r)=\sum_{l=0}^D r^{-l}.
\end{gathered}\tag{WCF14}
\]
Here $v!E_A(z)/(\mu_vz^v)$ differs from one by at most
$C_vB_0|z|e^{B_0|z|}/(v+1)$, by its exponential Taylor series.
This is at most $1/2$ for $|z|\le R_0$. On $|t|\le a_0$,
$|w(t)|\le R_0/2$ and $|w(t)/t|\ge d_0$, by the arctangent
series. Factoring $g$ as
$(\mu_v/v!)(w/t)^ve^{-4w}[v!E_A(w)/(\mu_vw^v)]$ proves
$|1/g(t)|\le M_0$ on that entire disk, including zero.

Fix any $0<\varrho<1$ with no zero of $g$ on $|t|=\varrho$.
Let its interior zeros be $\zeta$, of orders $m_\zeta$, and put
\[
\begin{gathered}
 H_\varrho=\sum_\zeta m_\zeta,\qquad
 A_{\zeta,h}=\frac1{(m_\zeta-h)!}
   \left.\frac{d^{m_\zeta-h}}{dt^{m_\zeta-h}}
       \frac{(t-\zeta)^{m_\zeta}}{g(t)}\right|_{t=\zeta},\\
 f_\varrho(t)=\frac1{g(t)}-
       \sum_{\zeta}\sum_{h=1}^{m_\zeta}\frac{A_{\zeta,h}}{(t-\zeta)^h},
 \quad W_\varrho=\operatorname{arctanh}\varrho,\\
 m_{E,\varrho}=\min_{|t|=\varrho}|E_A(w(t))|>0,\quad
 M_\varrho=\frac{\varrho^ve^{4W_\varrho}}{m_{E,\varrho}}
       +\sum_{\zeta,h}\frac{|A_{\zeta,h}|}{(\varrho-|\zeta|)^h}.
\end{gathered}\tag{WCF15}
\]
All these pole sums are finite. Zero is not a pole, and $f_\varrho$
is holomorphic on a neighbourhood of the closed circle and its
interior. Its modulus on the circle is at most $M_\varrho$.

Truncated power-series multiplication gives
$T_D(g)^{-1}=T_D(1/g)$. The inverse weighted upper diagonals have
ratios $\rho_{i+l}/\rho_{i+v}$. WCF11 gives an upper bound
$\sqrt{2l+1}$ at $s=1$. At $s=a$, this ratio is at most
$\rho_i/\rho_{i+v}\le\sqrt{(b)_v/v!}$; the last inequality follows
by decreasing each factor $(i+t+b)/(i+t+1)$ to its maximum at $i=0$.
Thus set
\[
 W_{D,s}^-=
 \begin{cases}\sqrt{2D+1},&s=1,\\\sqrt{(s/2)_v/v!},&s=a,
 \end{cases}
 \qquad B^0_{D,s}=W_{D,s}^-M_0S_D(a_0).
 \tag{WCF16}
\]
Cauchy's bound on $1/g$ at radius $a_0$, followed by the sum of
these upper-diagonal norms, proves $\|A_{D,s}^{-1}\|\le B^0_{D,s}$.

For completeness, the Taylor coefficient of $A/(t-\zeta)^h$ is
$c_n=A(-1)^h\zeta^{-n-h}\binom{n+h-1}{h-1}$.
Extend the binomial polynomial to every integer $n$. The full matrix
$(c_{j-i})_{0\le i,j\le D}$ has rank at most $h$, by expanding its
polynomial in $i$. For all orders at one pole the combined column
space is contained in $\operatorname{span}\{(\zeta^ii^d)_i:
0\le d<m_\zeta\}$, hence its rank is at most $m_\zeta$.
For a lower-diagonal gap $l=i-j\ge1$, the omitted entry equals
zero if $l<h$ and equals
$-A\zeta^{l-h}\binom{l-1}{h-1}$ if $l\ge h$.
Its cancellation therefore requires a lower-triangular correction with
the opposite, positive displayed coefficient.

After multiplying by the actual inverse diagonal factors, that
lower-diagonal shift has weight $\rho_j/\rho_{j+l+v}$. Its norm is
at most one for $s=1$, and at most
$\sqrt{(b)_{l+v}/(l+v)!}$ for $s=a$, by the same decreasing-factor
argument. Define the finite, original-pole quantity
\[
\begin{gathered}
 \Theta_{l,s}=\begin{cases}1,&s=1,\\
       \sqrt{(s/2)_{l+v}/(l+v)!},&s=a,\end{cases}\\
 Z_{D,s,\varrho}=\sum_{\zeta,h}|A_{\zeta,h}|
      \sum_{l=h}^{D}|\zeta|^{l-h}\binom{l-1}{h-1}\Theta_{l,s},\\
 B^{\rm rem}_{D,s,\varrho}=W_{D,s}^-M_\varrho S_D(\varrho)
                              +Z_{D,s,\varrho},\qquad
 B^{\rm cap}_{D,s,\varrho}
     =\min\{B^0_{D,s},B^{\rm rem}_{D,s,\varrho}\}.
\end{gathered}\tag{WCF17}
\]
An empty inner sum is zero. The preceding exact principal-part
calculation and the coefficient bound on $f_\varrho$ prove a decomposition
of the entire inverse into a matrix of rank at most $H_\varrho$ and
a remainder of norm at most $B^{\rm rem}_{D,s,\varrho}$.
Diagonal multiplication leaves this rank bound unchanged. No pole,
lower-triangular correction, or original coefficient has been removed.

For $c$ orthogonal copies, $c\in\{1,4\}$, and a rank $r$ plane or
attained subquotient, put $h_c=\min\{r,cH_\varrho\}$ and
\[
 \mathcal L^{\rm wt}_{D,s,r,c}(\varrho)
 =2h_c\log^+B^0_{D,s}
       +2(r-h_c)\log^+B^{\rm cap}_{D,s,\varrho},\qquad
 \mathcal U^{\rm wt}_{D,s,r}=2r\log^+B^+_{D,s}.
 \tag{WCF18}
\]
To prove the inverse exterior bound, take the kernel of the exceptional
matrix. It has codimension at most $cH_\varrho$, and on it the full
inverse has norm at most the remainder bound. Diagonalizing its positive
squared norm shows that every singular value after the first
$cH_\varrho$ is at most that bound; all singular values are at most
$B^0_{D,s}$. Multiplying the largest $r$ singular values proves the
exterior bound in WCF18, with $\log^+$ giving a nonnegative allowance.

The older inverse estimate remains available without alteration. In its
literal NCT16--19 notation define
\[
\begin{gathered}
 L_{s,D,v}=\tfrac12\max\{1,s/2-1\}\sum_{n=1}^{D+v}\frac1n,
 \quad G_{0,D}=M_0S_D(a_0),\\
 P_\varrho=\sum_{\zeta,h}\frac{|A_{\zeta,h}|}{(1-|\zeta|)^h},
 \quad V_{\varrho,D}=M_\varrho S_D(\varrho)+P_\varrho,\\
 \mathcal L^{\rm old}_{D,s,r,c}(\varrho)
 =4rL_{s,D,v}+2h_c\log^+G_{0,D}
       +2(r-h_c)\log^+\min\{G_{0,D},V_{\varrho,D}\},\\
 \mathcal L^{\rm pole}_{D,s,r,c}=
  \inf_{\varrho\ {\rm admissible}}
  \min\{\mathcal L^{\rm wt}_{D,s,r,c}(\varrho),
        \mathcal L^{\rm old}_{D,s,r,c}(\varrho)\}.
\end{gathered}\tag{WCF19}
\]
The old bound follows from the same pole decomposition before diagonal
weighting: summing the lower shifts gives $P_\varrho$ by the binomial
series, and each of the two diagonal factors costs at most
$e^{L_{s,D,v}}$. The latter follows by writing $\rho_n$ as its
product and using $\log(1+x)\le x$. Applying the exterior argument
then gives exactly WCF19. Thus taking the displayed minimum and
infimum keeps two independently proved finite bounds. Admissible radii
exist arbitrarily close to one because zeros are isolated. The infimum
is finite and nonnegative and need not be attained.

\subsection{The full determinant controls every inverse exterior rank}

There is a stronger estimate that uses the entire original triangular
matrix and needs no separation of its poles. Put $d=D+1$ and
$M_{D,s}=\max\{1,B^+_{D,s}\}$. Triangularity of WCF6 proves the
exact complex determinant identity
\[
 \det A_{D,s}=g(0)^d\prod_{i=0}^{D}\frac{\rho_{i+v}}{\rho_i}.
 \tag{WCF19a}
\]
It retains the complete phase $g(0)=(-i)^v\mu_v/v!$. For $s=1$ the
product of positive ratios is at least one. For $s=a$ its reciprocal
satisfies
\[
 \begin{split}
 \log\prod_{i=0}^{D}\frac{\rho_i}{\rho_{i+v}}
 &=\frac12\sum_{i=0}^{D}\sum_{t=0}^{v-1}
           \log\left(1+\frac{s/2-1}{i+t+1}\right)\\
 &\le\frac12(s/2-1)v\sum_{j=1}^{d}\frac1j.
 \end{split}\tag{WCF19b}
\]
For $v=0$ both sums are empty. Define the nonnegative finite quantity
\[
 J_{D,s}=d\log^+\frac1{|g(0)|}
       +\frac12(s/2-1)_+v\sum_{j=1}^{d}\frac1j.
 \tag{WCF19c}
\]
Equations WCF19a--b prove $|\det A_{D,s}|^{-1}\le e^{J_{D,s}}$.

Let $B=(A_{D,s}^{\mathsf T})^{\oplus c}$, of dimension $cd$, with
singular values $\sigma_1\ge\cdots\ge\sigma_{cd}>0$. Its inverse
exterior norm is
\[
 \left\|\bigwedge^rB^{-1}\right\|
 =\frac{\sigma_1\cdots\sigma_{cd-r}}
           {\sigma_1\cdots\sigma_{cd}}
 =\frac{\|\bigwedge^{cd-r}B\|}{|\det B|}
 \le M_{D,s}^{cd-r}e^{cJ_{D,s}}\qquad(0\le r\le cd).
 \tag{WCF19d}
\]
The first equality multiplies the $r$ largest inverse singular values;
the denominator is exactly the modulus of the full determinant. Thus
this estimate includes every direction of the original inverse, including
all directions caused by every symbol zero. It has imposed no zero-free
unit disk or bound on its pole count.

For $r>0$ set
\[
 \mathcal L^{\rm det}_{D,s,r,c}
 =2\bigl((cd-r)\log M_{D,s}+cJ_{D,s}\bigr),\qquad
 \mathcal L^*_{D,s,r,c}
 =\min\{\mathcal L^{\rm pole}_{D,s,r,c},
                         \mathcal L^{\rm det}_{D,s,r,c}\}.
 \tag{WCF19e}
\]
At $r=0$ define both allowances to be zero, as the actual empty
exterior norm and determinant are one. This sharpens rather than
replaces the preserved finite principal-part bounds. For fixed $C_{\rm cut}>0$, if
$D\le C_{\rm cut}a^2$, $c\le4$, $s\in\{1,a\}$, then WCF13 and WCF19c give
\[
 \mathcal L^*_{D,s,r,c}
 \le\mathcal L^{\rm det}_{D,s,r,c}
 =O_{\rm actual,C_{\rm cut}}(a^2\log(a+1))=o(aq)
 \quad(0\le r\le c(D+1)).
 \tag{WCF19f}
\]
Indeed $cd-r\le4(D+1)$, the first term is
$O(a^2\log(a+1))$, and $J_{D,s}=O_{\rm actual,C_{\rm cut}}(a^2+
a\log(a+1))$. The bounds hold uniformly over every retained rank.
No limiting choice of radius is needed for WCF19f.

\subsection{The original lower evaluation space and the conductor Gram}

Now take the original cutoffs $N\ge q-1$ and $D=N-v$. Define
the full observation matrices, with the indicated rows and fixed root order,
\[
 (\mathsf E_-)_{n,(u,w)}=\phi_{n,s,c_1}(\beta_{a-8,u,w}),\quad0\le n\le D,
 \qquad
 (\mathsf E_+)_{n,(u,w)}=\phi_{n,s,c_0}(\beta_{a,u,w}),\quad0\le n\le N.
 \tag{WCF20}
\]
The lower dimension is $q'$. Since $q-q'=16a-48\ge96>v$,
$D\ge q'-1$, so the lower matrix has full column rank by its
nonzero leading coefficients and the lower Vandermonde. The upper
matrix likewise has rank $q$. Let the full original conductor matrix
$C_a$ be defined by collecting every original shift:
\[
 \sum_{u,w}(C_a^{\mathsf T}\eta)_{uw}e^{\beta_{a,u,w}z}
       =E_A(z)\sum_{u,w=0}^{a-8}\eta_{uw}e^{\beta_{a-8,u,w}z}.
 \tag{WCF21}
\]
The identity $\beta_{a-8,u,w}+\beta_{8,r,t}=\beta_{a,u+r,w+t}$
proves the formula, retaining all terms with the same resulting root.
The map $C_a^{\mathsf T}$ is injective: a vanishing product makes
the lower exponential sum zero on an open set, hence everywhere,
and its distinct-root Vandermonde forces $\eta=0$.

Complex-linear functionals transform by ordinary transpose. In fact
evaluating $\mathcal T_A\phi_{v+j,s,c_0}$ at the lower roots and using
WCF6 gives exactly
\[
 \boxed{\quad \mathsf E_+ C_a^{\mathsf T}
       =\begin{bmatrix}0_{v\times q'}\\ A_{D,s}^{\mathsf T}\mathsf E_-\end{bmatrix},
 \qquad
 K_{C,N,s}=\overline C_a(\mathsf E_+^*\mathsf E_+)C_a^{\mathsf T}
       =\mathsf E_-^*\overline{A_{D,s}}A_{D,s}^{\mathsf T}\mathsf E_-.
 \quad}\tag{WCF22}
\]
The first $v$ rows vanish because those polynomials lie in the full
kernel of $\mathcal T_A$. This is the exact map from the original
lower evaluation space to the conductor-multiple subspace. Its lower
Gram is $K_{-,D,s}=\mathsf E_-^*\mathsf E_->0$ and its upper Gram
is the literal $K_{C,N,s}>0$ of TAC5. No new lower source order is
introduced: the two orders remain $1$ and $a$, even though the lower
root packet has degree $a-8$. The shift of its centre is exactly four.

For a full-column-rank matrix $E$ of rank $r$ and an invertible $B$,
the Gram determinant is the squared norm of the wedge of its columns.
Applying $\bigwedge^r B$ to this wedge and then applying
$\bigwedge^rB^{-1}$ to its image proves the two directed determinant
bounds. The transpose has the same singular values as its matrix:
if $A=U\Sigma V^*$ then $A^{\mathsf T}=\bar V\Sigma U^{\mathsf T}$.
Using WCF18--19 in WCF22 therefore proves
\[
 \boxed{\quad
 -\mathcal L_{D,s,q',1}^*
 \le e_{N,s}^{C}:=\log\det K_{C,N,s}-\log\det K_{-,N-v,s}
 \le\mathcal U^{\rm wt}_{D,s,q'}.
 \quad}\tag{WCF23}
\]
The same proof applies to restrictions and attained quotients. To see
the precise map, write $K\subset Y$ in the original space, identify
$Y/K$ with $Y\cap K^\perp$, and identify $BY/BK$ with
$BY\cap(BK)^\perp$. The induced map and its inverse are the
restrictions of $B$ and $B^{-1}$ followed by orthogonal projections.
Each projection contracts every exterior norm, which follows in an
orthonormal basis from its zero or unit singular values. Thus the
same bounds use the actual quotient rank, with the full original
denominator retained. The determinant/wedge identity is classical;
compare Lyons, \S4, (4.1) \cite{human:lyons2003}. The exact maps,
weights and their bounds have been proved here.

For any rank $r\le C_r a^2$, $D\le C_{\rm cut}a^2$, both original orders,
and $c=1,4$, WCF13 gives the new uniform forward conclusion
\[
 \mathcal U^{\rm wt}_{D,s,r}=O_{\rm actual,C_{\rm cut},C_r}(a^2\log(a+1))
                            =o(aq).
 \tag{WCF24}
\]
It applies in particular to the complete $q'$-dimensional conductor
restriction, as well as to all earlier rank-$O(a)$ native subquotients.

\subsection{The four signs and the remaining calculation}

Retain the four actual endpoints and the TAC dual orientation:
\[
 N_i=(q-1,q,2q-1,2q)_i,\quad D_i=N_i-v,\quad
 \mathcal R f=-f(N_0)-f(N_1)+f(N_2)+f(N_3).
 \tag{WCF25}
\]
Set $U_i^C=\mathcal U^{\rm wt}_{D_i,s,q'}$,
$L_i^C=\mathcal L_{D_i,s,q',1}^*$, and
$\mathcal C^-_{a,s}=\mathcal R\log\det K_{-,N-v,s}$.
Subtracting the two low endpoint inequalities WCF23 and adding the
two high endpoint inequalities gives the exact finite receiver
\[
 \boxed{\quad
 \mathcal C^-_{a,s}-U_0^C-U_1^C-L_2^C-L_3^C
 \le\mathcal C_{a,s}
 \le\mathcal C^-_{a,s}+L_0^C+L_1^C+U_2^C+U_3^C.
 \quad}\tag{WCF26}
\]
Here $\mathcal C_{a,s}=\mathcal R\log\det K_{C,N,s}$ is the exact
conductor term in TAC10 and TAC15. Every $U_i^C$ and $L_i^C$ is
$O_{\rm actual}(a^2\log(a+1))=o(aq)$, by WCF19f and WCF24.
Consequently the full conductor return transfers to the literal
lower-root covariance return with a proved two-sided error on that
scale. The receiver respects all four signs and every actual correlated
endpoint.

The finite endpoint dimensions give a sharper upper-directed statement.
At the two low endpoints $d_i=D_i+1$ satisfies
$d_0-q'=m$ and $d_1-q'=m+1$, where $m=16a-48-v$. Thus
\[
 \begin{split}
 \mathcal C_{a,s}-\mathcal C^-_{a,s}
 \le{}&2m\log M_{D_0,s}+2(m+1)\log M_{D_1,s}
            +2J_{D_0,s}+2J_{D_1,s}\\
       &+2q'\log M_{D_2,s}+2q'\log M_{D_3,s}.
 \end{split}\tag{WCF26a}
\]
This follows from WCF26 by using the determinant bound only on the
two negatively signed low endpoints. Its first line is
$O_{\rm actual}(q+a\log(a+1))$. Its second line is at most
$32\delta q\log a+O_{\rm actual}(q)$ for $s=a$, and at most
$(32\delta+4)q\log a+O_{\rm actual}(q)$ for $s=1$. To verify the
coefficients, use $D_2,D_3=2(a+1)^2-v+O(1)$ in WCF9 and WCF12:
$\log M_{D_i,a}\le8\delta\log a+O_{\rm actual}(1)$ and
$\log M_{D_i,1}\le(8\delta+1)\log a+O_{\rm actual}(1)$.
There are two upper terms, each multiplied by $2q'\le2q$.
These are upper bounds on an actual signed correction, not assigned
values of independent endpoint determinants.

For the analytic quotient of rank $m=q-q'-v=16a-48-v$, replace
$q'$ by $m$ in WCF18--19. Its new forward allowance is
$O(a\log(a+1))$. The retained old inverse estimates at rank $m$
have the already proved NCT21 radius-infimum limit $o(aq)$:
at any fixed admissible radius their limsup divided by $a^3$ is
bounded by a constant times $\log(1/\varrho)$, while the exceptional
rank is fixed; only afterwards let $\varrho$ approach one. Thus the
TAC10a--11 analytic correction may use the smaller forward constants
without changing its existing two-sided conclusion. The complete
conductor has rank $q'$ and is addressed instead by WCF26.

Finally let $F_{L_a}^{(s)}$, $\delta^{\rm ar}_{a,s}$, $\Xi_{a,s}$,
$\mathcal T^{\le R}_{a,s}$, $\mathscr T_{a,s,R}$, and $W_a,P_{a,\pm}$
be the unchanged terms defined by TAC9--15. Substitution into TAC15
gives the exact identity
\[
 \begin{split}
 J_a={}&\mathcal C^-_{a,s}+\mathcal R e^C_{N,s}
      +F_{L_a}^{(s)}+\delta^{\rm ar}_{a,s}+\Xi_{a,s}
      +\mathcal T^{\le R}_{a,s}+\mathscr T_{a,s,R}\\
      &+W_a-P_{a,-}-P_{a,+}.
 \end{split}\tag{WCF27}
\]
In particular the upper complete-action estimate receives
$L_0^C+L_1^C+U_2^C+U_3^C$ with coefficient one, and WCF26a retains
its more precise endpoint dimensions. The bound $o(aq)$ does not make
this correction negligible relative to TAC20's actual benchmark
$4q\log(\delta q(a+1)/2)=12q\log a+O_{\rm actual}(q)$.
The explicit $q\log a$ coefficients and the coupled actual terms must
therefore be retained. The next quantitative calculation is the
transferred lower-root covariance return together with the surviving
native correlated tail and arithmetic correction in WCF27; the generic
inverse transport estimate is now proved. WCF22 specifies its literal
lower space, source order, centre and shifted cutoffs. WCF15--17 retain
every pole and lower-triangular correction when finer signed control is
needed. The original independent
degree-$k$ proper-source matrices and all CPQ complements continue to
obey their existing exact receiving maps; this degree-$a$ calculation
has not substituted the target packet for that proper source.
\endgroup
